\chapter{Conclusións e Posibles Ampliacións}

\section{Grao de cumprimento dos obxectivos}

O proxecto acadou os obxectivos formulados na introdución de forma satisfactoria. Telaria centraliza nunha única
aplicación a xestión integral dunha viaxe en grupo, substituíndo o conxunto disperso de
ferramentas que se empregan habitualmente. O Cadro~\ref{tab:obxectivos} resume o grao de
cumprimento de cada obxectivo.

\begin{table}[H]
  \centering
  \caption{Grao de cumprimento dos obxectivos.}
  \label{tab:obxectivos}
  \begin{tabularx}{\textwidth}{@{}p{6.2cm}cX@{}}
    \toprule
    \textbf{Obxectivo} & \textbf{Estado} & \textbf{Observacións} \\
    \midrule
    Gastos compartidos e liquidacións & Cumprido & Cálculo no servidor con minimización de transferencias (RF10, RF11). \\[3pt]
    Calendario de eventos con localización & Cumprido & Eventos con data, duración e localización xeográfica (RF12). \\[3pt]
    Documentos compartidos & Cumprido & Subida directa ao almacén con miniaturas (RF13). \\[3pt]
    Chat de grupo & Cumprido & Mensaxería en tempo real mediante SSE (RF14). \\[3pt]
    Asistente de IA por viaxe & Cumprido & Chatbot con historial persistente sobre Ollama (RF15). \\[3pt]
    Rede de amizades (complementaria) & Cumprido & Solicitudes, lista de amigos e convite áxil (RF16, RF17). \\[3pt]
    Xestión de conta & Cumprido & Perfil, contrasinal, avatar e eliminación con confirmación (RF03, RF04). \\[3pt]
    Arquitectura API-first & Cumprido & Contrato OpenAPI con xeración de código en ambos os extremos. \\[3pt]
    Multiplataforma iOS/Android & Cumprido & Única base de código con React Native e Expo. \\
    \bottomrule
  \end{tabularx}
\end{table}

\section{Valoración técnica}

A adopción dunha arquitectura \textit{API-first} demostrou ser un acerto: ter o contrato
OpenAPI como fonte de verdade permitiu xerar automaticamente os tipos do cliente e as
interfaces do servidor, mantendo ambos os extremos sincronizados e detectando
incoherencias en tempo de compilación. O uso de \textit{Server-Sent Events} para o chat e
o asistente de IA, e o das URLs prefirmadas para os ficheiros, permitiron implementar
funcionalidades esixentes sen sobrecargar o servidor. A execución local do modelo de
linguaxe con Ollama evitou a dependencia de servizos externos de pago. Finalmente, a
batería de probas automatizadas do backend proporcionou unha rede de seguridade que
facilitou a evolución continua do código.

Entre as dificultades atopadas destacan a integración de Testcontainers cun entorno de
contedores sen daemon (Podman), a coordinación dos fluxos asíncronos de \textit{streaming}
e a xestión coherente da subida de ficheiros en varios pasos.

\section{Posibles ampliacións}

O sistema deixa aberto un conxunto de liñas de mellora:

\begin{itemize}
  \item \textbf{Notificacións push} para invitacións, mensaxes novas e recordatorios de eventos.
  \item \textbf{Preparación para produción}: externalizar os segredos de configuración
        (actualmente en ficheiros de configuración), habilitar HTTPS/TLS e facer
        configurable a dirección do servidor no cliente. Empregar servizos de hosting para o backend que permitan escalar a demanda,
        mellores modelos para o asistente de IA, monitorización de métricas e alertas, etc. O frontend
            tamén debería ser empaquetado e publicado en \textit{App Store} e \textit{Google Play}.
  \item \textbf{Engadir as probas do frontend ao pipeline de desenvolvemento}: Debido a limitacións 
       de tempo e á tardía implementación das probas do frontend polo seu carácter máis propenso a cambios e
       rediseños, non se puideron integrar estas no pipeline de desenvolvemento.
   A súa integración permitiría detectar erros de interface e de interacción de usuario antes de que cheguen á produción.
  \item \textbf{Modo sen conexión} con sincronización posterior, para mellorar a
        experiencia en destinos con conectividade limitada.
  \item \textbf{Accesibilidade}: Mellorar a accesibilidade tendo en conta diversos perfís de
        usuario con distintas dificultades (engadir un modo para daltónicos, etc.).
  \item \textbf{Integración de provedores de identidade}: Permitir o inicio de sesión mediante provedores externos como Google ou Meta.
  \item \textbf{Soporte web}: Aínda que a aplicación está pensada para dispositivos móbiles, un cliente web permitiría acceder á información dende calquera navegador.
\end{itemize}
